% =====================================================================
%  Chương 3 — Mục 3 của đề: hai tập ảnh và một tập về bệnh tiểu đường
%  Mã nguồn: a05/data.py · Notebook: notebooks/02_du_lieu.ipynb
% =====================================================================
\chapter{Dữ liệu}
\label{ch:dulieu}

Mục 3 của đề yêu cầu hai tập ảnh và một tập về bệnh tiểu đường, có đường dẫn và mô tả. Bảng~\ref{tab:batap}
tóm tắt ba tập được dùng. Hai đường dẫn đã được kiểm tra còn truy cập được (HTTP 200) ngày 25/9/2026.

\begin{table}[H]
  \centering
  \caption{Ba tập dữ liệu của bài}
  \label{tab:batap}
  \small
  \begin{tabular}{L{2.3cm}L{6.6cm}L{5.8cm}}
    \toprule
    Tập & Nguồn & Quy mô \\
    \midrule
    CIFAR-10 & \url{https://www.cs.toronto.edu/~kriz/cifar.html}~\cite{krizhevsky2009} &
      \SL{data/cifar10/n_train}\,+\,\SL{data/cifar10/n_val}\,+\,\SL{data/cifar10/n_test} ảnh $32\times32\times3$,
      \SL{data/cifar10/n_classes} lớp \\
    CIFAR-100 & cùng trang trên, tệp \texttt{cifar-100-python.tar.gz} &
      \SL{data/cifar100/n_train}\,+\,\SL{data/cifar100/n_val}\,+\,\SL{data/cifar100/n_test} ảnh $32\times32\times3$,
      \SL{data/cifar100/n_classes} lớp mịn, 20 siêu lớp \\
    Diabetes & \url{https://www.kaggle.com/datasets/iammustafatz/diabetes-prediction-dataset}~\cite{mustafa2023} &
      \SL{data/diabetes/n_raw} dòng, \SL{data/diabetes/n_features_raw} đặc trưng, nhãn nhị phân \\
    \bottomrule
  \end{tabular}
\end{table}

\begin{phathien}{Về tên ``MNIST-100'' trong đề}
Đề ghi hai tập ảnh là CIFAR-10 và ``MNIST-100''. Không có bộ dữ liệu chuẩn nào mang tên MNIST-100. Bài hiểu
đây là \textbf{CIFAR-100}: cùng tác giả, cùng kích thước ảnh và cùng định dạng tệp với CIFAR-10, chỉ khác số lớp
(100 thay vì 10). Cặp CIFAR-10/CIFAR-100 cũng cho một phép thử có ý nghĩa: cùng kiến trúc, cùng cấu hình, chỉ tăng
số lớp lên mười lần.
\end{phathien}

\section{CIFAR-10}

CIFAR-10~\cite{krizhevsky2009} gồm 60.000 ảnh màu tự nhiên kích thước $32\times32$ thuộc 10 lớp. Mỗi lớp có đúng
\SL{data/cifar10/per_class_max} ảnh huấn luyện và \SL{data/cifar10/test_per_class} ảnh kiểm tra, tức hoàn toàn cân bằng.
Hình~\ref{fig:cifar10mau} là ảnh huấn luyện đầu tiên của mỗi lớp, vẽ lại từng điểm ảnh bằng pgfplots từ tệp dữ liệu mà
notebook 02 xuất ra.

\begin{figure}[H]
  \centering
  \hinh{f3_cifar10_mau}
  \caption{Mỗi lớp CIFAR-10 một ảnh, vẽ từ $32\times32$ giá trị RGB}
  \label{fig:cifar10mau}
\end{figure}

\section{CIFAR-100}

CIFAR-100 có cùng 60.000 ảnh $32\times32$ nhưng chia thành \SL{data/cifar100/n_classes} lớp mịn, mỗi lớp chỉ
\SL{data/cifar100/per_class_max} ảnh huấn luyện, ít hơn CIFAR-10 mười lần. 100 lớp mịn được gom thành 20 siêu lớp,
mỗi siêu lớp đúng \SL{data/cifar100/fine_per_coarse} lớp mịn. Mô hình được huấn luyện trên nhãn mịn; siêu lớp chỉ
dùng ở Chương~\ref{ch:sosanh} để vẽ ma trận nhầm lẫn ở mức 20 nhóm cho dễ đọc.

\begin{figure}[H]
  \centering
  \hinh{f3_cifar100_mau}
  \caption{Mỗi siêu lớp CIFAR-100 một ảnh, kèm tên siêu lớp}
  \label{fig:cifar100mau}
\end{figure}

\section{Mô tả và tiền xử lý hai tập ảnh}

\paragraph{Chia tập.} 50.000 ảnh huấn luyện được tách \SL{data/cifar10/n_train} ảnh train và \SL{data/cifar10/n_val}
ảnh validation, phân tầng theo nhãn, hạt giống 42. Validation dùng để chọn checkpoint. \SL{data/cifar10/n_test} ảnh test chính thức
chỉ được dùng một lần, ở cuối.

\maNguon{data__split_train_val}{Tách validation phân tầng}

\paragraph{Chuẩn hoá.} Mỗi kênh được trừ trung bình rồi chia độ lệch chuẩn, cả hai tính trên nhánh train. Với CIFAR-10,
trung bình ba kênh R, G, B là \SL{data/cifar10/mean/0}, \SL{data/cifar10/mean/1}, \SL{data/cifar10/mean/2} và độ lệch chuẩn
là \SL{data/cifar10/std/0}, \SL{data/cifar10/std/1}, \SL{data/cifar10/std/2} (theo thang $[0,1]$). Con số của CIFAR-100 lần lượt
là \SL{data/cifar100/mean/0}, \SL{data/cifar100/mean/1}, \SL{data/cifar100/mean/2} và \SL{data/cifar100/std/0},
\SL{data/cifar100/std/1}, \SL{data/cifar100/std/2}. Hình~\ref{fig:histpixel} cho thấy kênh xanh dương lệch về phía tối
hơn hai kênh còn lại ở cả hai tập, và có đỉnh ở giá trị 255 do nền trời và nền trắng.

\begin{figure}[H]
  \centering
  \hinh{f3_hist_pixel}
  \caption{Phân bố cường độ điểm ảnh theo kênh trên nhánh train}
  \label{fig:histpixel}
\end{figure}

\paragraph{Nạp dữ liệu.} Hai tệp \texttt{.tar.gz} gốc được đọc thẳng bằng \texttt{tarfile} và \texttt{pickle},
không giải nén ra đĩa, rồi lưu đệm thành \texttt{.npz}.

\maNguon{data__load_cifar}{Đọc CIFAR-10/100 từ tệp gốc}

\paragraph{Dữ liệu nằm sẵn trên GPU.} Trên Windows, \texttt{DataLoader} nhiều tiến trình vừa chậm vừa dễ lỗi. Ở dạng
\texttt{uint8}, cả tập ảnh đủ nhỏ để nằm gọn trong VRAM, nên được đưa hết lên GPU một lần. Mỗi lô được cắt ra bằng chỉ số
ngẫu nhiên, chuẩn hoá và tăng cường ngay trên GPU. Tăng cường gồm hai phép: đệm 4 điểm ảnh rồi cắt ngẫu nhiên lại
$32\times32$ (dịch ảnh tối đa 4 điểm theo mỗi chiều), và lật ngang với xác suất 0,5. Chỉ nhánh train được tăng cường.

\maNguon{data__gpu_augment}{Tăng cường dữ liệu chạy trên GPU}
\maNguon{data__GPUImageData.batches}{Sinh lô từ tensor nằm sẵn trên GPU}

\paragraph{Dạng tensor.} Mỗi lô đưa vào mạng có dạng \texttt{[B, 3, 32, 32]} kiểu \texttt{float32}, $B=256$.

\section{Diabetes Prediction Dataset}

\subsection{Mô tả}

Tập Diabetes Prediction~\cite{mustafa2023} gồm \SL{data/diabetes/n_raw} hồ sơ bệnh nhân, mỗi hồ sơ có
\SL{data/diabetes/n_features_raw} đặc trưng và nhãn \texttt{diabetes} (1 là mắc bệnh):

\begin{itemize}
  \item Nhân khẩu: \texttt{gender} (\SL{data/diabetes/categories/gender/0}, \SL{data/diabetes/categories/gender/1},
        \SL{data/diabetes/categories/gender/2}) và \texttt{age}.
  \item Tiền sử: \texttt{hypertension}, \texttt{heart\_disease} (0/1) và \texttt{smoking\_history} (6 mức).
  \item Chỉ số cơ thể và xét nghiệm: \texttt{bmi}, \texttt{HbA1c\_level} (\%), \texttt{blood\_glucose\_level} (mg/dL).
\end{itemize}

Tập có \SL{data/diabetes/n_dup} dòng trùng lặp hoàn toàn. Các dòng này bị bỏ trước khi chia tập, nếu không một bản
sao có thể nằm ở train còn bản kia nằm ở test. Sau khi bỏ trùng còn \SL{data/diabetes/n_rows} dòng, trong đó
\SL{data/diabetes/n_pos} dòng dương tính (\SL{data/diabetes/pos_rate}\,\%). Mất cân bằng này quyết định cách đánh giá
ở Chương~\ref{ch:sosanh}: trên tập test, đoán mọi người đều khoẻ đã đạt accuracy \SL{diabetes/majority_acc}\,\%, nên accuracy
một mình không nói được gì.

\begin{table}[H]
  \centering
  \caption{Thống kê mô tả bốn đặc trưng số (sau khi bỏ trùng)}
  \label{tab:diabdesc}
  \input{generated/tab_diab_desc}
\end{table}

Giá trị \texttt{"No Info"} của \texttt{smoking\_history} chiếm \SL{data/diabetes/no_info_rate}\,\% số dòng. Đây không phải
giá trị thiếu ngẫu nhiên mà là một nhóm bệnh nhân không khai báo, nên được giữ lại như một mức riêng.

\subsection{Quan hệ giữa đặc trưng và nhãn}

Hình~\ref{fig:diabhist} tách phân bố của ba đặc trưng theo nhãn. Hai chỉ số xét nghiệm tách hai lớp rõ nhất. Notebook 02
tìm được hai ngưỡng mà phía trên đó \emph{mọi} bệnh nhân đều mắc bệnh: HbA1c từ \SL{data/diabetes/pure_threshold/HbA1c_level}\,\%
và đường huyết từ \SL{data/diabetes/pure_threshold/blood_glucose_level} mg/dL. Có \SL{data/diabetes/pos_covered}\,\% ca dương
tính vượt ít nhất một trong hai ngưỡng này, nên phần lớn ca dương rất dễ nhận ra; phần khó của bài toán là số ca còn lại,
có chỉ số xét nghiệm nằm trong vùng chồng lấn với người khoẻ. Histogram của hai chỉ số này có dạng răng cưa vì trong
tập dữ liệu chúng chỉ nhận một số ít giá trị rời rạc, không phải số đo liên tục. Ma trận tương quan (Hình~\ref{fig:diabcorr}) cho kết quả khớp: hệ số tương quan với nhãn cao
nhất là của \texttt{blood\_glucose\_level} (\SL{data/diabetes/corr_target/blood_glucose_level}) và \texttt{HbA1c\_level}
(\SL{data/diabetes/corr_target/HbA1c_level}), còn \texttt{age} là \SL{data/diabetes/corr_target/age}.

\begin{figure}[H]
  \centering
  \hinh{f3_diabetes_hist}
  \caption{Phân bố HbA1c, đường huyết và tuổi theo hai lớp (mật độ, sau khi bỏ trùng)}
  \label{fig:diabhist}
\end{figure}

\begin{figure}[H]
  \centering
  \hinh{f3_diabetes_corr}
  \caption{Hệ số tương quan Pearson giữa các đặc trưng số, nhị phân và nhãn}
  \label{fig:diabcorr}
\end{figure}

\subsection{Tiền xử lý và dạng tensor}

Sau khi bỏ trùng, dữ liệu được chia 70/15/15 phân tầng theo nhãn thành \SL{data/diabetes/n_train} dòng train,
\SL{data/diabetes/n_val} dòng validation và \SL{data/diabetes/n_test} dòng test. Mọi phép biến đổi có tham số đều
\emph{chỉ học trên nhánh train} rồi áp dụng nguyên cho hai nhánh còn lại: \texttt{StandardScaler} cho bốn cột số, và danh
mục one-hot cho \texttt{gender} và \texttt{smoking\_history}. Kết quả là \SL{data/diabetes/n_features_encoded} đặc trưng.

\maNguon{data__prepare_diabetes}{Chuẩn bị dữ liệu Diabetes, không rò rỉ từ test}

Mỗi bệnh nhân là một vectơ $\mathbf{x}\in\mathbb{R}^{\SL{data/diabetes/n_features_encoded}}$, được đưa vào mạng dưới dạng
\texttt{[B, 1, \SL{data/diabetes/n_features_encoded}]}: một kênh, chiều dài bằng số đặc trưng. Bộ lọc Conv1D kích thước 3
trượt dọc các cột đặc trưng kề nhau. Thứ tự cột là cố định (số, nhị phân, rồi one-hot) nhưng không mang ý nghĩa không gian
như điểm ảnh. Chương~\ref{ch:sosanh} sẽ cho thấy điều này ảnh hưởng thế nào đến lợi ích của các cơ chế kiến trúc.
